\section{OmniTrace}
\label{sec:method}

\subsection{Problem Setting}
\label{sec:problem}

\paragraph{Input and source units.}
A decoder-only omni-modal LLM receives an input sequence $\mathbf{x}=(x_1,\dots,x_N)$ in which text tokens and the encoder outputs for images, audio, and video share one sequence.
We split the attributable part of $\mathbf{x}$ into \emph{source units} $\mathcal{S}=\{S_1,\dots,S_m\}$, each a set of input positions, using structure the input already has.
A text passage is one unit and an image is one unit.
Audio and video are cut into units of about one second, following the time grid of the model's own encoder, so each unit $S_j$ has a time interval $[b_j,e_j]$.
The system prompt, the question, the answer options, and template tokens are part of the prompt and are never candidates.
Each unit belongs to one modality, and we write $\mathcal{M}(S_j)$ for the set of units that share the modality of $S_j$.

\paragraph{Output.}
The model generates $\mathbf{y}=(y_1,\dots,y_T)$ with greedy decoding. We split $\mathbf{y}$ into sentences $C_1,\dots,C_K$, each a set of generated-token indices.
For every sentence $C_k$, and for the response as a whole, we want (i) a score for every unit that ranks the units by how much the model relied on them, and (ii) a small labelled set $\hat{\mathcal{S}}_k\subseteq\mathcal{S}$, reported as source ids for text and images and as time intervals for audio and video.
A sentence may get an empty set.

\subsection{Reading Attention as Mass over Units}
\label{sec:scoring}

\paragraph{Principle.}
When the model predicts $y_t$, each attention head distributes one unit of weight over the positions it reads from.
Let $A^{(\ell,h)}_t(i)$ be the weight on input position $i$ in head $h$ of decoder layer $\ell$.
We average over all heads and over the last third of the layers, $\mathcal{L}_{\text{late}}$:
\begin{equation}
a_t(i)=\frac{1}{|\mathcal{L}_{\text{late}}|\,H}\sum_{\ell\in\mathcal{L}_{\text{late}}}\sum_{h=1}^{H}A^{(\ell,h)}_t(i),
\end{equation}
so $a_t$ is a distribution over where the token $y_t$ reads from.
A source unit is a set of positions, and the natural question is how much of this distribution falls on the set.
OmniTrace scores a unit by this \emph{attention mass}, summed over the tokens of a sentence:
\begin{equation}
s_k(S_j)=\sum_{t\in C_k}\;\sum_{i\in S_j} a_t(i).
\label{eq:mass}
\end{equation}
The score can be read as the expected number of the sentence's tokens that read from $S_j$.
It is additive in two ways. Over tokens, the response-level score is the sum of the sentence scores, so one pass gives attributions at both levels. Over positions, the score of two merged units is the sum of their scores, so the scores do not depend on how a stream is cut.
Nothing else enters: no token weights, no normalization over sources, and no fitted parameter.

\paragraph{Mass, not peaks.}
\label{sec:design}
A common way to read attention is to trace each generated token to its single highest-weight position and to count these traces as votes.
For text sources the peak and the mass tell the same story, because a source is a few tokens.
For continuous sources they do not. One second of video is spread over about 300 positions, each of which receives little weight; the second as a whole can hold much of the mass while no position in it wins the argmax, which goes to a sharply attended text token.
We measured this on 1{,}000 development questions of NExT-GQA that share no video with our test set.
Ranking video seconds by argmax votes agrees with human-marked evidence at chance level in every layer (AUROC 0.50--0.52), where mass reaches 0.60; and when the argmax ranges over the whole prompt, the question text wins about 98\% of the tokens.

\paragraph{Late layers.}
On the same data, the first twelve layers of Qwen2.5-Omni carry no source signal (AUROC 0.50--0.53) and averaging all layers dilutes the late ones (0.57 against 0.60).
We therefore fix the band as the last third of the depth, which is defined for any model (layers 18--27 of 28 in Qwen2.5-Omni-7B, 24--35 of 36 in MiniCPM-o~4.5), and do not tune it further.

\paragraph{Cost.}
Scores are read during generation. At each step we keep only the attention row of the token being generated in the late layers, $O(|\mathcal{L}_{\text{late}}|\,N)$ memory per step.
OmniTrace adds no forward pass, no backward pass, and no training to the generation it explains.

\subsection{From Scores to Evidence Labels}
\label{sec:labels}

Scores rank the units. To give the user a short answer, we label the units that stand out from their competitors.
Let $\bar{s}_k^{\,\setminus j}$ be the mean score of the other units of the same modality as $S_j$, that is, the mean of $s_k(S')$ over $S'\in\mathcal{M}(S_j)\setminus\{S_j\}$.
A unit is labelled for sentence $C_k$ when its score is at least $\tau$ times this mean:
\begin{equation}
S_j\in\hat{\mathcal{S}}_k \iff s_k(S_j)\;\ge\;\tau\,\bar{s}_k^{\,\setminus j}.
\label{eq:label}
\end{equation}
Three properties follow from this rule.
First, units compete only within their modality, so text passages cannot crowd out images, whose scores are on a different scale.
Second, the number of labelled units adapts to the sentence: flat scores give an empty label, so a sentence that draws on nothing in particular abstains.
Third, the rule is relative, so it does not depend on the absolute scale of attention, which changes with input length and model.
Adjacent labelled time units are merged into intervals $[b,e]$; text and image units are listed by id.
We use $\tau=1.5$ for every modality and both models (\Cref{appendix:threshold}).
The sum in \Cref{eq:mass} favors large units when units differ much in size; we measure this in \Cref{sec:size}.
